\section{Conclusion and Discussion}\label{sec:discussion}

In this work, we demonstrate poking manipulation as a fundamental motion primitive that complements grasping and encompasses pushing in terms of capability. Our work is the first to show qualitative and quantitative results for multiple test conditions to demonstrate the flexibility and robustness of poking as a skill through \textsl{PokeRRT}. We present the task times, number of executed actions, and success rates of our proposed motion planners and four baseline algorithms across six different scenarios. The results demonstrate the strengths of the poking motion primitive: poking is not as limited by robot reachability, robot end-effector design, object properties, and inaccuracies due to perception as grasping or pushing. Success rates are higher for poking-based planners than for the push planner indicating that poking expands the size of reachable workspace by its ability to execute longer object displacements using shorter robot end-effector trajectories. Task times for computed plans are significantly lower for poking than for pushing, indicating that poking allows fast object manipulation because it does not face the same constant-contact restriction as pushing. \textsl{PokeRRT} and \textsl{PokeRRT*} also demonstrated the qualitative behaviors visualized in \cref{fig:scenarios}.

Errors in dynamic modeling caused by the gap between simulation and the real-world are covered by the replanning strategy in our work. However, future work in identifying sources of uncertainty will allow for the quantification of object--obstacle collision margins and incorporating them directly into the planning process will lower the chances of collision in the real-world. In order to make plan execution in the real-world more robust, our future work will also formalize analytical and learned models for poking and quantify object constraints which will lead to additional insights about the capabilities and limitations of poking. For instance, poking may not be an effective mode of manipulation for high friction interactions and should also be used sparingly since it may cause significant wear-and-tear on a real robot through repeated execution of high-speed trajectories. Consequently, by discovering the feasibility conditions of multiple skills (e.g. poking, pushing, and grasping) through interaction and planning around the key strengths of each motion primitive, robots will be able to more efficiently manipulate objects. 
\vspace{-2pt}
